\section{Monitoring framework and evaluation design}
\label{sec:framework}

This section defines what both studies share: the context policies, the
per-turn quality signals, the streaming monitor and the evidence standards
(Figure~\ref{fig:framework}).

\begin{figure}[!htbp]
\centering
\begin{tikzpicture}[
  font=\small, >=Latex, node distance=4mm and 6mm,
  box/.style={draw, rounded corners=2pt, align=center, minimum height=9mm,
              inner sep=3pt, fill=#1},
  box/.default=white]
\node[box=gray!10, text width=27mm] (user) {User turn $u_t$\\(new information,\\question)};
\node[box=blue!8, text width=33mm, right=of user] (policy)
  {Context policy $\pi$\\[1pt]\footnotesize dialogue history $\cdot$ fresh prefix\\
   \footnotesize $\cdot$ verified state summary};
\node[box=gray!10, text width=24mm, right=of policy] (llm) {LLM\\primary reply\\+ $K{-}1$ samples\\+ $L$ rewordings};
\node[box=orange!10, text width=26mm, below=9mm of llm] (metrics)
  {Per-turn quality\\vector $\mathbf{q}_t$\\\footnotesize CQ CS RS LS ES SS};
\node[box=orange!10, text width=27mm, left=of metrics] (window)
  {Window means\\$\bar q_{m,t}$ over the last\\$W$ valid values};
\node[box=red!8, text width=30mm, left=of window] (detect)
  {Drift detectors\\\footnotesize CUSUM, PH, EWMA,\\\footnotesize DDM, ADWIN};
\node[box=green!10, text width=30mm, below=8mm of detect] (alarm)
  {Alarm time $\tau_m$\\\footnotesize compared with the change $c$};
\node[box=green!10, text width=27mm, right=of alarm] (null)
  {Threshold from\\stationary nulls\\\footnotesize false alarms $\le\alpha$};
\node[box=gray!5, text width=26mm, right=of null] (gold)
  {Exact gold $y_t$\\\footnotesize(evaluation only)};
\draw[->] (user) -- (policy);
\draw[->] (policy) -- (llm);
\draw[->] (llm) -- (metrics);
\draw[->] (metrics) -- (window);
\draw[->] (window) -- (detect);
\draw[->] (detect) -- (alarm);
\draw[->, dashed] (null) -- (alarm);
\draw[->, dashed] (gold) -- (null);
\draw[->, dashed] (gold.north) -- (metrics.south);
\draw[->, dotted] (llm.north) -- ++(0,4mm) -| node[pos=0.25, above, font=\footnotesize]
  {primary reply enters later history (dialogue history only)} (policy.north);
\end{tikzpicture}
\caption{The monitoring framework. A context policy decides what the model
sees at turn $t$. Responses are scored into a per-turn quality vector, smoothed
over a window and passed to drift detectors whose thresholds are calibrated on
stationary dialogues. Exact gold is used only to score accuracy and evaluate
alarms, never by the monitor.}
\label{fig:framework}
\end{figure}

\subsection{Context policies}

At turn $t$ the user sends $u_t$ and the model returns a primary reply $r_t$
from a context $C_t$ set by a policy $\pi$:
\begin{align}
  \text{dialogue history:}\quad & C_t = (u_1, r_1, \dots, r_{t-1}, u_t),\\
  \text{fresh prefix:}\quad & C_t = (u_1, \dots, u_t),\\
  \text{verified state summary:}\quad & C_t = (s_{t-1}, u_t),
\end{align}
where $s_{t-1}$ is the exact task state after turn $t-1$, rendered by a
solver. Dialogue history versus fresh prefix estimates the \emph{net} effect of
retaining assistant-side history; context length and message positions differ
as well, so the effect is not attributed to the content of the replies alone.
The verified state summary shortens the context, removes earlier replies and
supplies a correct state at once; it tests mitigation, not mechanism. Our code
packs every context into the request, with server-side memory, retrieval and
prompt rewriting disabled, after a pilot in which a gateway silently dropped
native message history. Every context is therefore reproducible from the
stored records.

\subsection{Per-turn quality signals}

Each primary reply is accompanied by $K-1$ further samples from the same
context and by $L$ replies to frozen rewordings of the instruction. RQEval
\citep{rqeval2026} computes six dimensions in $[0,1]$, used unchanged:
accuracy (CQ, the primary answer equals the gold); consistency (CS, pairwise
agreement of the $K$ answers); robustness (RS, agreement of reworded replies
with a correct primary answer); logical coherence (LS, absence of NLI
contradictions between consecutive steps; \citealp{he2021deberta}); efficiency
(ES, correctness combined with token economy); and stability (SS, BERTScore
similarity of the $K$ traces; \citealp{zhang2020bertscore}). CQ, RS and ES
require a final answer and gold, CS needs at least two answers, and LS and SS
need only reasoning text. This matters for early warning: a signal that
requires an answer is unobservable before the model is ready to give one.
Missing values are never imputed and are reported with their support.

\subsection{The streaming monitor}

For each metric $m$, the monitor keeps a trailing mean $\bar q_{m,t}$ over the
last $W\in\{10,20,30\}$ valid values ($W=20$ primary); missing values neither
enter nor advance the window. Five detectors, CUSUM, Page--Hinkley (PH), EWMA,
DDM and ADWIN, are paired with metrics by three strategies taken from the
funded project design (Table~\ref{tab:strategies}). Thresholds are never tuned
on drift data: for each metric, detector and window, we take the most
sensitive grid value whose first-alarm rate on stationary streams does not
exceed $\alpha$. Study~1 uses permutation nulls with family-wise
$\alpha=0.05$ verified on disjoint permutations. Study~2 uses its stationary
dialogues, with leave-one-out false-alarm estimates. For a change at turn $c$,
an alarm before $c$ is a false alarm and an alarm in $[c, c+33]$ a detection.
A signal gives \emph{early warning} only if it detects drifts in time, respects
the null calibration and does so across models.

\begin{table}[!htbp]
\centering
\caption{Detector assignment strategies. All monitor a decrease of the window
mean; CQ is additionally monitored on its raw binary stream by DDM.}
\label{tab:strategies}
\footnotesize
\setlength{\tabcolsep}{4pt}
\begin{tabular}{lcccccc}
\toprule
Strategy & CQ & CS & RS & LS & ES & SS\\
\midrule
Distribution-based & CUSUM & CUSUM & CUSUM & CUSUM & CUSUM & CUSUM\\
Proposal mapping & CUSUM & ADWIN & DDM & PH & EWMA & ADWIN\\
ADWIN everywhere & ADWIN & ADWIN & ADWIN & ADWIN & ADWIN & ADWIN\\
\bottomrule
\end{tabular}
\end{table}

\subsection{A paired readiness measure}

Sharded tasks allow an answer-free paired signal. Let $a_t^{\mathrm{F}}$ and
$a_t^{\mathrm{D}}$ be 1 if the fresh-prefix and dialogue-history arms attempt an
answer at turn $t$, and 0 if they state that more information is needed. The
\emph{Cumulative Resolution Lag} of a task is
\begin{equation}
  \mathrm{CRL} = \frac{1}{|\mathcal{T}|}\sum_{t\in\mathcal{T}}
  \max\!\left(a_t^{\mathrm{F}}-a_t^{\mathrm{D}},\,0\right),
\end{equation}
over the pre-final turns $\mathcal{T}$ at which both arms responded validly.
CRL measures how persistently the history arm lags an information-matched
counterfactual in becoming ready to answer. It was frozen as an auxiliary
measure before the Study~1 data were collected and is not an RQEval dimension.

\subsection{Evidence standards}

Every result carries one of four labels. \emph{Confirmatory} results test
hypotheses frozen before data collection; \emph{prespecified secondary}
results are frozen but outside the confirmatory family; \emph{exploratory}
results use rules fixed before computation on data already seen; and
\emph{post-hoc} results were designed after inspection. Protocols fixed the
sample, prompts, models, scoring, missing-data policy and tests, and were
hashed together with all input records. Raw responses are stored before scoring
and never overwritten, failed gates are reported rather than relabelled, and
every number, table and figure is regenerated by scripts from the hashed
inputs.
